% ============================================================
% 50_typography_semantics.tex -- Schriftmakros und semantische Helfer
% ============================================================

\newcommand{\ZSFfontChapter}{\sffamily\bfseries\fontsize{8.1bp}{9.0bp}\selectfont}
\newcommand{\ZSFfontSection}{\sffamily\bfseries\normalsize}
% Dritte Ebene = kompakter, aber klar sichtbarer Lookup-Locator.
\newcommand{\ZSFfontSubsubsection}{\sffamily\bfseries\fontsize{7.8bp}{8.6bp}\selectfont}
\newcommand{\ZSFfontBoxTitle}{\sffamily\bfseries\normalsize}
\newcommand{\ZSFfontNote}{\sffamily\normalsize\itshape}
\newcommand{\ZSFfontCodeComment}{\ttfamily\color{ce_green}}
\newcommand{\ZSFfontTitleHeader}{\sffamily\bfseries\fontsize{10.5bp}{11.8bp}\selectfont}
\newcommand{\ZSFfontTitleSubheader}{\sffamily\itshape\normalsize}
\newcommand{\ZSFfontTitleTag}{\sffamily\itshape\normalsize}
\newcommand{\ZSFfontDangerUpper}{\sffamily\bfseries}

\newcommand{\TableCellNote}[1]{{\ZSFfontNote #1}}
\newcommand{\TableTitleBreak}{\newline}
\newcommand{\TableTitle}[2][]{%
  \textbf{#2}%
  \if\relax\detokenize{#1}\relax\else\TableTitleBreak\TableCellNote{#1}\fi
}
\renewcommand{\ZSFfontTable}{\normalsize}
\renewcommand{\ZSFfontTableDense}{\normalsize}

% Index-sichere Basis-API: ohne Modul bleiben alle Varianten reine
% Hervorhebung bzw. unsichtbare No-ops. styles/66_index.tex ersetzt diese
% Definitionen nur bei aktivem Stichwortverzeichnis.
\NewDocumentCommand{\ZSFkeyword}{s o m}{\textbf{#3}}
\NewDocumentCommand{\ZSFindex}{o m}{}
\NewDocumentCommand{\ZSFindexsee}{o m m}{}
\newcommand{\ZSFScriptRef}[1]{\hfill\textcolor{\ChapterBarTextColor!80}{\normalsize (S.#1)}}
\newcommand{\ZSFhl}[1]{\underline{\textbf{#1}}}

\makeatletter
\def\ZSF@refExtractFirst#1#2\@nil{#1}
\def\ZSF@refExtractChap#1.#2\@nil{#1}
\newcommand{\ZSF@refComputeTargetColor}[1]{%
  \@ifundefined{r@#1}{%
    \edef\ZSF@refTargetColor{\chaptercolor}%
  }{%
    \edef\ZSF@refTargetNum{\expandafter\expandafter\expandafter
      \ZSF@refExtractFirst\csname r@#1\endcsname\@nil}%
    \edef\ZSF@refTargetChap{\expandafter\ZSF@refExtractChap\ZSF@refTargetNum.\@nil}%
    \ZSFSetPaletteColorNameFromNumber{\ZSF@refTargetChap}{\ZSF@refTargetColor}%
  }%
}
\newcommand{\ZSFref}[1]{%
  \begingroup
  \ZSF@refComputeTargetColor{#1}%
  \hyperref[#1]{{\ZSFfontNote\upshape\bfseries\color{\ZSF@refTargetColor}($\rightarrow$\,\ref*{#1})}}%
  \endgroup
}
\newcommand{\ZSFsectionref}[1]{%
  \begingroup
  \ZSF@refComputeTargetColor{#1}%
  \hyperref[#1]{{\upshape\bfseries\color{\ZSF@refTargetColor}\ref*{#1}}}%
  \endgroup
}
\newcommand{\ZSFarrowref}[1]{%
  \begingroup
  \ZSF@refComputeTargetColor{#1}%
  \hyperref[#1]{{\upshape\bfseries\color{\ZSF@refTargetColor}(\ensuremath{\rightarrow}~\ref*{#1})}}%
  \endgroup
}
\makeatother

\newcommand{\ZSFconclusion}[1]{%
  \par\vspace{-\parskip}\vspace{\ZSFspaceXS}%
  \noindent$\Rightarrow$ #1%
  \par\vspace{-\parskip}\vspace{\ZSFspaceXS}%
}

\newcommand{\ZSFrowColor}{\rowcolor{\FormulaboxBackColor}}
\newcommand{\ZSFtableRuleColor}{black!70}
\newcommand{\SetZSFtableRuleColor}[1]{\renewcommand{\ZSFtableRuleColor}{#1}}
\newcommand{\ZSFtableRuleWidth}{0.5pt}
\newcommand{\SetZSFtableRuleWidth}[1]{\renewcommand{\ZSFtableRuleWidth}{#1}}
\newcommand{\ZSFbeginTableRuleStyle}{%
  \begingroup
  \arrayrulecolor{\ZSFtableRuleColor}%
  \setlength{\arrayrulewidth}{\ZSFtableRuleWidth}%
}
\newcommand{\ZSFendTableRuleStyle}{\endgroup}

\makeatletter
\renewcommand{\section}{\@ifstar{\zsf@section@star}{\zsf@section@plain}}
\newcommand{\zsf@section@plain}[1]{\StartChapter{#1}}
\newcommand{\zsf@section@star}[1]{\StartFrontChapter{#1}}

\renewcommand{\subsection}{\@ifstar{\zsf@subsection@star}{\zsf@subsection@plain}}
\newcommand{\zsf@subsection@plain}[1]{\SubsectionBar{#1}}
\newcommand{\zsf@subsection@star}[1]{\SubsectionBar*{#1}}

\renewcommand{\subsubsection}{\@ifstar{\zsf@subsubsection@star}{\zsf@subsubsection@plain}}
\newcommand{\zsf@subsubsection@plain}[1]{\SubsubsectionBar{#1}}
\newcommand{\zsf@subsubsection@star}[1]{\SubsubsectionBar*{#1}}
\makeatother
